\cotefile{1}
\scanpage{122}
\[
n!\alpha_n\leq\frac{\varepsilon_0}{n+1}\,\varepsilon_1\ldots\varepsilon_n\,e^n
\qquad \sqrt[n]{n!\alpha_n}\leq e\left(\frac{\varepsilon_0}{n+1}\right)^{1/n}
\deleted{\unsure{\varepsilon_1\ldots\varepsilon_n}}\unsure{(\varepsilon_1\ldots\varepsilon_n)^{1/n}|(\varepsilon_1\ldots\varepsilon_n)},
\]
\unsure{(} $n\to\infty$, $\left(\frac{\varepsilon_0}{n+1}\right)^{1/n}\to1$, et
$(\varepsilon_1\ldots\varepsilon_n)^{1/n}\to0$,
comme on voit immédiatement, d'où $\sqrt[n]{n!\alpha_n}\to0$,
d'où suit que $\unsure{S(\lambda)=}\sum n!\alpha_n(\lambda)$ \unsure{converge absolument} \unsure{à} \unsure{ch.} \unsure{en résulte|en vertu}
% slanted
\margin{par suite \unsure{convergence|converge} à \unsure{distance} \unsure{finie} sur $l^1$, la fonction $S(\lambda)$ est holomorphe}
Comme on a $|\alpha_n(\lambda)|\leq\|\lambda\|_1^n$, il suit que
la fonction $S(\lambda)$ \deleted{\unsure{sur $l^1$}} est holomorphe au voisinage
de l'origine\,; d'ailleurs, la forme \unsure{$n$ fois linéaire}
$\alpha_n(\lambda^{(1)},\ldots,\lambda^{(n)})$ associée à $\alpha_n(\lambda)$ est
\[
\alpha_n(\lambda^{(1)}\ldots\lambda^{(n)})=\frac{1}{n!}
\sum_{i_1,\ldots,i_n}\lambda^{(1)}_{i_1}\ldots\lambda^{(n)}_{i_n}
\]
(la somme étant étendue aux systèmes $(i_1,\ldots,i_n)$ d'indices \unsure{deux à deux|tous} distincts),
d'où $\|\alpha_n\|\leq\frac{1}{n!}$ \inserted{\deleted{\unsure{$\|\alpha_n\|$}} $n!\|\alpha_n\|\leq1$},
le rayon de convergence strict $R_0$ de la fonction $S(\lambda)$ à l'origine est
donc \deleted{\unsure{être|à peu}} $\geq1$. Soit alors $a_{\deleted{\unsure{m}}}^{(m)}$ l'élément de $l^1$
\deleted{\unsure{égal à}} donné par
\[
a_i^{(m)}=\begin{cases}\frac{1}{m}&\text{si }i\leq m\\0&\text{si }i>m\end{cases}
\]
on a manifestement $\|a^{(m)}\|_1=1$, et
\[
\alpha_n(a^m)=\deleted{\unsure{\frac{1}{m^n}}}
\left\{\begin{aligned}
&\deleted{\unsure{0\quad\text{si }n>m}}\\
&\frac{1}{m^n}\frac{m!}{n!(m-n)!}\quad\text{si }n\leq m
\end{aligned}\right.\,,\quad\text{d'où}
\]
\[
n!\alpha_n(a^m)=\frac{1}{m^n}\,m(m-1)\ldots(m-n+1)=
\deleted{\unsure{1}}(1-\tfrac{1}{m})(1-\tfrac{2}{m})\ldots(1-\tfrac{n-1}{m})
\]
Donc $n!\alpha_n(a^{(m)})\to1$ pour $m\to\infty$.
Il suit que $S(a^{(m)})\to\infty$ si $m\to\infty$,
donc $S(\lambda)$ n'est pas borné dans la boule
unité de $l^1$, et son rayon de convergence
large $R_1$ à l'origine est donc aussi $\leq1$, d'où $R_0=R_1=1$.

\deleted{Enfin, $S(\lambda)$ étant holomorphe au voisinage de \unsure{$0$|$a$} et pouvant être \unsure{prolongée|holomorphe|holomorphes} sur \unsure{chaque droite}}
% archive sheet number: 123
